\section{Stratification générationnelle, supports de quarks et réalisation neutrino}
\label{sec:generaciones-quarks-neutrinos-rev2}
\index{générations fermioniques!stratification HMT--MD}
\index{quark!supports cellulaires et registre}
\index{PMNS!transport des bases}

La loi générale agit sur des multisections et non sur une liste de masses.
L'apparition de trois générations s'exprime donc comme compatibilité de
trois structures : la stratification de l'atlas, les registres de saveur et
l'espace tridimensionnel de mélange. Dans le secteur leptonique chargé, les
multisections
\begin{equation}
 G_0^{\ell},\qquad G_2^{\ell},\qquad G_4^{\ell}
 \label{eq:tres-multisecciones-leptonicas-rev2}
\end{equation}
ont pour rangs \(1,8,16\). Le centre \(G_0^\ell\) contient la section
électronique de rang un ; \(G_2^\ell\) et \(G_4^\ell\) conservent, en revanche,
des spectres opératoriels non triviaux. L'organisation générationnelle complète
est
\begin{center}
\small
\begin{tabularx}{0.96\textwidth}{@{}c c c c X@{}}
\toprule
\textbf{Génération} & \textbf{Lepton} & \textbf{Par quark} &
\textbf{Sabor neutro} & \textbf{Soporte estructural}\\
\midrule
1 & \(e\) & \((u,d)\) & \(\nu_e\) & centre \(G_0^\ell\), registre de
première génération et premier vecteur de la base d'interaction\\
2 & \(\mu\) & \((c,s)\) & \(\nu_\mu\) & multisection \(G_2^\ell\),
registre de deuxième génération et deuxième vecteur de la base d'interaction\\
3 & \(\tau\) & \((t,b)\) & \(\nu_\tau\) & multisection \(G_4^\ell\),
registre de troisième génération et troisième vecteur de la base d'interaction\\
\bottomrule
\end{tabularx}
\end{center}
La conjugaison particule--antiparticule inverse les coordonnées de charge et
d'orientation de la représentation, mais conserve cette stratification ; par
conséquent, elle ne double pas le nombre de générations.

La dimension trois dispose en outre d'un porteur d'incidence propre. La
représentation orientée dodécaphasique contient
\begin{equation}
 V_{12}\cong V_1\oplus V_3\oplus V_{3'}\oplus V_5,
 \qquad \operatorname{rank}P_3=3,
 \label{eq:descomposicion-rango-tres-generaciones-rev2}
\end{equation}
et l'entrelaceur dodécaphasique--Witt transporte le secteur positif selon
\begin{equation}
 P_G\xleftrightarrow{\;W_{GK}\;}P_3
 \xleftrightarrow{\;U_W\;}Q_3,
 \qquad
 \operatorname{rank}P_G=\operatorname{rank}P_3
 =\operatorname{rank}Q_3=3.
 \label{eq:transporte-rango-tres-generaciones-rev2}
\end{equation}
Ce porteur tridimensionnel permet de comparer les bases générationnelles. Les
chemins, les signatures et les opérateurs de masse fournissent ensuite le contenu
spectral que CKM et PMNS transportent.

\subsection{Position interne et registres des \(6\) saveurs de quark}

La position d'une saveur de quark est la donnée composée
\begin{equation}
 \mathfrak q_X=
 (\text{branche de charge},\text{génération},(\mathbb Z/3\mathbb Z)_{\rm color},
  \text{multisection},\text{registre de route}).
 \label{eq:posicion-quark-compuesta-rev2}
\end{equation}
Les branches de charge \(-1/3\) possèdent les supports
\begin{align}
 D_1&=\{(4,5),(6,5)\},\nonumber\\
 D_2&=\{(4,3),(4,7),(6,3),(6,7)\},\nonumber\\
 D_3&=\{(2,3),(2,5),(2,7),(8,3),(8,5),(8,7)\},\nonumber\\
 D_4&=\{(2,1),(2,9),(4,1),(4,9),(6,1),(6,9),(8,1),(8,9)\},
 \label{eq:soportes-quark-down-rev2}
\end{align}
tandis que les branches de charge \(+2/3\) occupent
\begin{align}
 U_1&=\{(4,4),(4,6),(6,4),(6,6)\},\nonumber\\
 U_3&=\{(2,2),(2,4),(2,6),(2,8),(4,2),(4,8),
       (6,2),(6,8),(8,2),(8,4),(8,6),(8,8)\}.
 \label{eq:soportes-quark-up-rev2}
\end{align}
Chaque support se relève à quatre orientations par cellule. Le quotient de couleur
organise \(36=12_{\rm R}+12_{\rm G}+12_{\rm B}\) biographies orientées et conserve une même
masse par saveur avant la réalisation de couleur.

Le transducteur nominal
\begin{equation}
 \mathcal N_{\rm ruta}(\Gamma_f)
 =((N,E,S,O),h_{369},r_5)
 \in\mathbb Z_{\ge0}^{4}\times\mathbb Z_{\ge0}^{2}
 \label{eq:transductor-nominal-quark-rev2}
\end{equation}
publie les coordonnées suivantes avant toute réalisation scalaire :
\begin{center}
\small
\begin{tabularx}{0.96\textwidth}{@{}c c c c c X@{}}
\toprule
\textbf{Saveur} & \textbf{Génération} & \textbf{Branche} &
\textbf{\((N,E,S,O)\)} & \textbf{\((h_{369},r_5)\)} &
\textbf{Eje dominante}\\
\midrule
\(u\)&1&\(+2/3\)&\((34,31,31,12)\)&\((51,13)\)&N--S\\
\(d\)&1&\(-1/3\)&\((27,35,18,28)\)&\((47,17)\)&E--O\\
\(c\)&2&\(+2/3\)&\((20,38,24,26)\)&\((49,16)\)&E--O\\
\(s\)&2&\(-1/3\)&\((50,32,12,14)\)&\((57,11)\)&N--S\\
\(t\)&3&\(+2/3\)&\((46,27,16,19)\)&\((57,9)\)&N--S\\
\(b\)&3&\(-1/3\)&\((12,63,9,24)\)&\((47,7)\)&E--O\\
\bottomrule
\end{tabularx}
\end{center}
Les signatures nominales associées sont
\begin{center}
\small
\begin{tabular}{@{}c c c@{}}
\toprule
\textbf{Saveur} & \textbf{\((n_A,s_C)\)} & \textbf{\((W_+,W_-)\)}\\
\midrule
\(u\)&\((11,7)\)&\((73,59)\)\\
\(d\)&\((17,8)\)&\((110,94)\)\\
\(c\)&\((61,8)\)&\((374,358)\)\\
\(s\)&\((41,-3)\)&\((243,249)\)\\
\(t\)&\((100,-1)\)&\((599,601)\)\\
\(b\)&\((71,-5)\)&\((421,431)\)\\
\bottomrule
\end{tabular}
\end{center}
Ces coordonnées distinguent des saveurs qui partagent une branche ou une génération et closent
le morphisme de saveur sans consulter aucune masse. Pour une section physique (P),
nous définissons
\begin{equation}
 \begin{aligned}
 \operatorname{sabor}(P)
 &=\bigl(\sigma_{\rm ud}(P),g(P),\chi_{\rm fina}(P)\bigr),\\
 \chi_{\rm fina}(P)
 &=(\text{rotations},\text{commutations},\text{axe},N,E,S,O,h_{369},r_5),
 \end{aligned}
 \label{eq:morfismo-sabor-ruta-cerrado-rev2}
\end{equation}
et, pour les quarks, on conserve en outre l'orbite ternaire
\([\gamma_P]_{\rm col}=\{\gamma_P^R,\gamma_P^G,\gamma_P^B\}\). La réalisation
physique est l'application déjà déterminée par les données précédentes,
\begin{equation}
 \mathcal S_{\rm phys}(P,\mathfrak M):
 \bigl(\sigma_{\rm ud},g,\chi_{\rm fina},[\gamma_P]_{\rm col}\bigr)
 \longmapsto \gamma^{\rm surv}_{P,\mathfrak M}
 \longmapsto
 \bigl(\sigma(\gamma^{\rm surv}_{P,\mathfrak M}),
       B_F^D,B_F^\Omega\bigr),
 \label{eq:realizacion-fisica-sabor-ruta-cerrada-rev2}
\end{equation}
où \(\mathfrak M\) est le couplage compatible et
\(\gamma^{\rm surv}_{P,\mathfrak M}\) le chemin qui persiste sous celui-ci. Pour
les leptons chargés, l'orbite de couleur se réduit à l'objet terminal. Le caractère
de masse s'évalue alors sur cette même généalogie :
\begin{equation}
 \boxed{
 \frac{m_{\mathfrak M}(P)}{m_e^{\rm HMT}}
 =\chi_{\rm mass}\!\left(
   \sigma\!\left(\gamma^{\rm surv}_{P,\mathfrak M}\right)\right).}
 \label{eq:masa-como-caracter-persistencia-rev2}
\end{equation}
Par conséquent, saveur, génération et couleur n'attendent pas un sélecteur futur : elles publient la
classe de chemin et son vecteur propre avant d'ouvrir la comparaison externe. La masse ne
choisit pas rétrospectivement le chemin ; elle ne fait que lire multiplicativement la persistance
que le couplage a stabilisée.

\subsection{Fermionicité, valeurs propres de masse et transport PMNS}

La clôture extérieure du foncteur d'échange construit les relations CAR
et fixe la fermionicité. Dans le secteur neutre, l'étiquette
\(\nu_e,\nu_\mu,\nu_\tau\) désigne une base d'interaction ; les valeurs de
masse sont des valeurs propres d'un opérateur neutre mélangé. Les référentiels internes
\(F_\ell,F_\nu:\mathbb C^3\to\operatorname{im}P_3\) comparent ces bases, et
la matrice
\begin{equation}
 U_{\rm HMT}=F_\ell^*\Phi_LF_\nu
 \label{eq:transporte-pmns-generaciones-rev2}
\end{equation}
est un transport unitaire. PMNS change de base entre les états d'interaction
et les vecteurs propres de masse ; l'échelle neutre appartient au spectre de l'opérateur,
non au changement de base.

\subsection{Extracteur neutre Z0 antérieur à la projection}
\label{sec:extractor-neutro-z0-torsion-rev2}

La construction récupérée du secteur neutre agit sur le registre
dodécaphasique enrichi, avant la contraction du feuillet, de la torsion et du bit
bidirectionnel. Soient \(E_m\), \(m\in\mathbb Z/12\mathbb Z\), les douze
événements de neuf ticks du registre CT--108. À chaque tick \(t\), on conserve
\[
 \theta(t)=\mathbf1_{\{d(t)=9\}},\qquad
 \varepsilon(t)\in\{-1,+1\},\qquad
 \kappa(t)=\mathbf1_{\{9\ {\rm en\ corona}\}},
 \qquad \omega(t)=(-1)^{\kappa(t)}.
\]
Le premier facteur détecte la torsion nonadique, le deuxième conserve le feuillet
two--ways et le troisième change le signe lors de la rotation de \(180^\circ\) de la
couronne. L'entier torsionnel de chaque événement et sa différence antipodale sont
\begin{equation}
 T_\nu(m)=\sum_{t\in E_m}
   \omega(t)\varepsilon(t)\theta(t),\qquad
 e_m^\nu=T_\nu(m)-T_\nu(m+6),
 \label{eq:extractor-z0-torsion-antipodal-rev2}
\end{equation}
et le trit neutre est fixé, sans paramètres continus, par
\begin{equation}
 \tau_m^\nu=\operatorname{sgn}(e_m^\nu)
 \in\{-1,0,+1\}.
 \label{eq:trit-neutro-torsion-antipodal-rev2}
\end{equation}

Pour les quatre orbites de saut quatre
\[
 I_a=\{a,a+4,a+8\}\pmod{12},\qquad a=0,1,2,3,
\]
le coefficient d'événement qui entre dans la signature et l'observable de majorité
des signes sont, respectivement,
\begin{equation}
 b_{120}^\nu
 =\sum_{a=0}^{3}\operatorname{sgn}
   \left(\sum_{m\in I_a}e_m^\nu\right),
 \qquad
 \nu_{120}^{\nu,\mathrm{sgn}}
 =\sum_{a=0}^{3}\operatorname{sgn}
   \left(\sum_{m\in I_a}\tau_m^\nu\right).
 \label{eq:canales-120-neutros-distintos-rev2}
\end{equation}
Ces deux entiers ne sont pas identifiés : le premier est la coordonnée
\(b_{120}=\nu_{120}^{\rm ev}\) du caractère, tandis que le second est un
diagnostic de majorité. La mémoire quadratique est
\begin{equation}
 b_\zeta^\nu=\nu_{270}^\nu
 =\sum_{m=0}^{11}\tau_m^\nu\tau_{m+3}^\nu.
 \label{eq:canal-270-neutro-rev2}
\end{equation}
Par conséquent, pour chaque histoire neutre survivante
\(\gamma_{\nu_i}\), le registre publie avant toute comparaison la signature
\[
 v_\nu(\gamma_{\nu_i})
 =\bigl(0,0,k_{\Delta_4}^{(i)},b_{120}^{\nu_i},
        b_\zeta^{\nu_i}\bigr)
\]
et l'exposant relatif
\begin{equation}
 \mathcal A_\nu^{(i)}
 =k_{\Delta_4}^{(i)}\Delta_4
  +b_{120}^{\nu_i}s_{120}
  +b_\zeta^{\nu_i}\zeta_{270},
 \qquad
 \zeta_{270}=135\Delta_4^2.
 \label{eq:caracter-relativo-neutro-z0-rev2}
\end{equation}
Ici, \(\zeta_{270}\) n'est pas remplacé par le moment de tour
\(s_{270}=q_-^{270}-q_+^{270}\).

La conjugaison globale two--ways
\(\tau^\nu\mapsto-\tau^\nu\) inverse
\(b_{120}^\nu\) et conserve \(b_\zeta^\nu\). Elle engendre donc deux
branches de spectre relatif,
\begin{equation}
 \mathcal A_{\nu,\pm}^{(i)}
 =k_{\Delta_4}^{(i)}\Delta_4
  \pm b_{120}^{\nu_i}s_{120}
  +b_\zeta^{\nu_i}\zeta_{270},
 \qquad
 \text{NON}\longleftrightarrow\text{IO},
 \label{eq:bifurcacion-neutra-no-io-rev2}
\end{equation}
qui se distinguent par le signe ultérieur de
\(\Delta m_{31}^2\), non par une sélection avec des données PDG. Cette construction
clôt l'extracteur discret Z0 et la bifurcation relative. L'échelle absolue
commune \(m_0\) et le noyau physique PMNS de rang complet restent des objets
séparés : ils ne sont pas introduits dans le sélecteur et ne sont pas déclarés
déduits de la seule parité two--ways.

La réalisation avec une unique phase de rang un produit
\[
 (\theta_{12},\theta_{23},\theta_{13})
 =(19.6807^\circ,56.4970^\circ,29.6224^\circ)
\]
et ne reproduit pas la phénoménologie du mélange. Ce contrôle sélectionne comme domaine approprié un
noyau de registre de rang complet,
\begin{equation}
 G_{ab}=\frac1{\sqrt{|L_a||N_b|}}
 \sum_{c\in L_a}\sum_{d\in N_b}
 K_{\rm registro}(c,d;\xi_c,\xi_d),
 \qquad U_{\rm int}=\operatorname{polar}(G),
 \label{eq:kernel-registro-rango-completo-rev2}
\end{equation}
dont la source technique est développée dans le chapitre du transport CKM et
PMNS. Ainsi sont séparés l'opérateur de masse, ses valeurs propres et l'application de
lecture par saveur.
